\documentclass[answers,12pt,addpoints]{exam}
\usepackage{import}
\import{C:/Users/prana/OneDrive/Desktop/MathNotes}{style.tex}
\usepackage{amsmath,amssymb,booktabs,longtable,pdflscape}

% Header
\newcommand{\name}{Pranav Tikkawar}
\newcommand{\course}{01:XXX:XXX}
\newcommand{\assignment}{Homework 3}
\author{\name}
\title{\course\ - \assignment}
\date{Due Thursday, October 22, 2026}

\begin{document}
\maketitle

\noindent Assigned exercises: 1.1, 1.2, and 2.1--2.5 from
\emph{A Course in Regression and Smoothing Methods}.

\medskip
\noindent Computational note: the numerical reference values below are those
printed in the exercise, not independently executed results. The accompanying
\texttt{homework3.R} implements all computations. Running it generates tables
that this document automatically includes when they are available.

\newpage
\begin{questions}

\question Exercise 1.1. For $x,a\in\mathbb R^p$ and a constant matrix
$A\in\mathbb R^{p\times p}$, find the gradients of $a^Tx$ and $x^TAx$.
\begin{solution}
\textbf{(i) Linear form.}
Since $a^Tx=\sum_{j=1}^p a_jx_j$, for each $k$,
\[
\frac{\partial}{\partial x_k}(a^Tx)=a_k.
\]
Stacking the partial derivatives as a column vector gives
\[
\boxed{\frac{\partial}{\partial x}(a^Tx)=a.}
\]

\textbf{(ii) Quadratic form.}
Expand the scalar function as
\[
x^TAx=\sum_{j=1}^p\sum_{k=1}^p A_{jk}x_jx_k.
\]
For $\ell\in\{1,\ldots,p\}$, the product rule gives
\[
\frac{\partial}{\partial x_\ell}(x^TAx)
=\sum_{k=1}^p A_{\ell k}x_k+\sum_{j=1}^p A_{j\ell}x_j
=(Ax)_\ell+(A^Tx)_\ell.
\]
Consequently,
\[
\boxed{\frac{\partial}{\partial x}(x^TAx)=(A+A^T)x.}
\]
If $A=A^T$, this reduces to $2Ax$.

For an alternative derivation, use differentials:
\begin{align*}
d(x^TAx)&=(dx)^TAx+x^TA\,dx\\
&=\bigl[(A+A^T)x\bigr]^Tdx.
\end{align*}
Comparison with $df=(\nabla f)^Tdx$ gives the same gradient.
Only the symmetric part $A_s=(A+A^T)/2$ contributes to the
quadratic form, since $x^T(A-A^T)x=0$.
\end{solution}

\question Exercise 1.2. Show that the two matrix-derivative expressions
for the Hessian of a scalar function agree, and find the Hessian of $x^TAx$.
\begin{solution}
\textbf{(i) Equality of the Hessian expressions.}
Assume $h$ has continuous second partial derivatives in a neighborhood of $x$.
Write
\[
L=\frac{\partial}{\partial x^T}\left(\frac{\partial h}{\partial x}\right),
\qquad
R=\frac{\partial}{\partial x}\left(\frac{\partial h}{\partial x^T}\right).
\]
By the derivative conventions in the exercise,
\[
L_{jk}=\frac{\partial}{\partial x_k}
\left(\frac{\partial h}{\partial x_j}\right),
\qquad
R_{jk}=\frac{\partial}{\partial x_j}
\left(\frac{\partial h}{\partial x_k}\right).
\]
Equality of mixed partials gives $L_{jk}=R_{jk}$ for every $j,k$.
Therefore,
\[
\boxed{
\frac{\partial}{\partial x^T}\left(\frac{\partial h}{\partial x}\right)
=\frac{\partial}{\partial x}\left(\frac{\partial h}{\partial x^T}\right)
=\nabla^2h(x).}
\]
In particular, $\nabla^2h(x)$ is symmetric. Merely assuming existence of
second partials at a point would not, by itself, justify interchanging them;
continuous second partials are a sufficient condition.

\textbf{(ii) Quadratic-form Hessian.}
By Exercise 1.1,
\[
\nabla(x^TAx)=(A+A^T)x.
\]
Differentiating this linear vector-valued function yields
\[
\boxed{\nabla^2(x^TAx)=A+A^T.}
\]
If $A$ is symmetric, the Hessian is $2A$.
\end{solution}

\question Exercise 2.1. Find the minimum-variance unbiased estimator
of the form $T_a=(1-a)Z_1+aZ_2$, where the independent variables satisfy
$Z_j\sim N(\mu,\sigma_j^2)$ and $\sigma_j^2>0$.
\begin{solution}
For every fixed $a$,
\[
E(T_a)=(1-a)\mu+a\mu=\mu.
\]
By independence,
\[
V(a)=\operatorname{Var}(T_a)
=(1-a)^2\sigma_1^2+a^2\sigma_2^2.
\]
Differentiating,
\[
V'(a)=-2(1-a)\sigma_1^2+2a\sigma_2^2.
\]
The first-order condition is
\[
(1-a)\sigma_1^2=a\sigma_2^2,
\]
so
\[
\boxed{a^*=\frac{\sigma_1^2}{\sigma_1^2+\sigma_2^2},\qquad
1-a^*=\frac{\sigma_2^2}{\sigma_1^2+\sigma_2^2}.}
\]
Since $V''(a)=2(\sigma_1^2+\sigma_2^2)>0$, this is the unique global minimum.
Equivalently,
\[
\boxed{T^*=\frac{\sigma_1^{-2}Z_1+\sigma_2^{-2}Z_2}
{\sigma_1^{-2}+\sigma_2^{-2}}.}
\]
Substitution gives
\[
\boxed{\operatorname{Var}(T^*)
=\frac{\sigma_1^2\sigma_2^2}{\sigma_1^2+\sigma_2^2}
=\frac{1}{\sigma_1^{-2}+\sigma_2^{-2}}.}
\]

\emph{Correction to the printed statement.}
For the estimator $(1-a)Z_1+aZ_2$, the correct proportionality is
\[
(1-a^*,a^*)\propto(\sigma_1^{-2},\sigma_2^{-2}),
\]
not $(a^*,1-a^*)\propto(\sigma_1^{-2},\sigma_2^{-2})$.
The latter would correspond to the estimator $aZ_1+(1-a)Z_2$.

If $\sigma_j^2=c_j\sigma^2$, then
\[
a^*=\frac{c_1}{c_1+c_2},\qquad 1-a^*=\frac{c_2}{c_1+c_2},
\]
so the common unknown scale cancels. Normality is unnecessary for this
variance-minimization argument; the common mean, independence, and finite
positive variances suffice.
\end{solution}

\question Exercise 2.2. Derive Newton--Raphson and Fisher scoring
for probit regression.
\begin{solution}
Let $Y_i\mid x_i\sim\operatorname{Bernoulli}(p_i)$ independently, where
\[
\eta_i=x_i^T\beta,\qquad p_i=\Phi(\eta_i),\qquad
q_i=1-p_i,\qquad f_i=\phi(\eta_i).
\]
Here $\Phi$ and $\phi$ are the standard normal CDF and density, and
$X$ has rows $x_i^T$.

\textbf{Step 1: Log-likelihood and score.}
The log-likelihood is
\[
\ell(\beta)=\sum_{i=1}^n\{y_i\log p_i+(1-y_i)\log q_i\}.
\]
Since $dp_i/d\eta_i=f_i$,
\[
u_i:=\frac{\partial\ell_i}{\partial\eta_i}
=\frac{y_if_i}{p_i}-\frac{(1-y_i)f_i}{q_i}
=\frac{f_i(y_i-p_i)}{p_iq_i}.
\]
Thus
\[
\boxed{U(\beta)=\nabla\ell(\beta)=\sum_i x_iu_i=X^Tu.}
\]

\textbf{Step 2: Observed information.}
Set $s_i=2y_i-1$, $t_i=s_i\eta_i$, and
$r(t)=\phi(t)/\Phi(t)$. Then $\ell_i=\log\Phi(t_i)$ and
\[
r'(t)=-r(t)\{t+r(t)\},\qquad u_i=s_ir(t_i).
\]
It follows that
\[
w_i^{\mathrm{obs}}=-\frac{\partial^2\ell_i}{\partial\eta_i^2}
=r(t_i)\{r(t_i)+t_i\}.
\]
Equivalently,
\[
w_i^{\mathrm{obs}}=
\begin{cases}
f_i^2/p_i^2+\eta_if_i/p_i,&y_i=1,\\
f_i^2/q_i^2-\eta_if_i/q_i,&y_i=0.
\end{cases}
\]
Writing $W_{\mathrm{obs}}=\operatorname{diag}(w_i^{\mathrm{obs}})$,
\[
J(\beta)=-\nabla^2\ell(\beta)=X^TW_{\mathrm{obs}}X.
\]
A first-order expansion of the score around $\beta^{(k)}$ gives
\[
0\approx U(\beta^{(k)})-J(\beta^{(k)})
(\beta^{(k+1)}-\beta^{(k)}).
\]
Therefore the Newton--Raphson update is
\[
\boxed{\beta^{(k+1)}=\beta^{(k)}+
(X^TW_{\mathrm{obs}}^{(k)}X)^{-1}X^Tu^{(k)}.}
\]

\textbf{Step 3: Fisher information and scoring.}
Under the model, $E(u_i\mid x_i)=0$ and
\[
E(u_i^2\mid x_i)
=\frac{f_i^2}{p_i^2q_i^2}\operatorname{Var}(Y_i\mid x_i)
=\frac{f_i^2}{p_iq_i}.
\]
Thus, with $w_i^F=f_i^2/(p_iq_i)$,
\[
I(\beta)=X^TW_FX,\qquad W_F=\operatorname{diag}(w_i^F).
\]
Replacing observed information by expected information gives
\[
\boxed{\beta^{(k+1)}=\beta^{(k)}+
(X^TW_F^{(k)}X)^{-1}X^Tu^{(k)}.}
\]

\textbf{Step 4: IRLS representation.}
Define the working response
\[
z_i^{(k)}=\eta_i^{(k)}+\frac{y_i-p_i^{(k)}}{f_i^{(k)}}.
\]
Since $w_i^F(z_i-\eta_i)=u_i$, scoring is equivalent to
\[
\boxed{\beta^{(k+1)}=(X^TW_F^{(k)}X)^{-1}X^TW_F^{(k)}z^{(k)}.}
\]
For probit regression, observed and expected information generally differ,
so the two algorithms follow different paths but target the same finite MLE.
In implementation, linear systems are solved instead of forming inverses.
\end{solution}

\question Exercise 2.3. Implement logistic Newton--Raphson, logistic MM,
probit Newton--Raphson, and probit Fisher scoring on \texttt{birthwt};
compare the prescribed starts and stopping rules.
\begin{solution}
\textbf{Step 1: Data and design.}
Use the binary response \texttt{low} and the formula
\[
\texttt{low}\sim\texttt{age+lwt+race+smoke+ptd+ht+ui+ftv}.
\]
Race has reference level white; $\texttt{ptd}=1\{\texttt{ptl}>0\}$;
and \texttt{ftv} is grouped into $0$, $1$, and $2+$, with reference $0$.
This gives $n=189$ observations and $p=11$ coefficients, including the intercept.
All four algorithms use the same design matrix.

\textbf{Step 2: Logistic Newton--Raphson.}
For $p_i=\operatorname{expit}(x_i^T\beta)$,
\[
\ell(\beta)=\sum_i\{y_ix_i^T\beta-\log(1+e^{x_i^T\beta})\},
\quad U=X^T(y-p),\quad J=X^TWX,
\]
where $W=\operatorname{diag}\{p_i(1-p_i)\}$. Hence
\[
\boxed{\beta^{(k+1)}=\beta^{(k)}+
(X^TW^{(k)}X)^{-1}X^T(y-p^{(k)}).}
\]

\textbf{Step 3: Logistic quadratic MM.}
Let $L=-\ell$. Because $p_i(1-p_i)\leq1/4$,
\[
\nabla^2L(\beta)\preceq\tfrac14 X^TX.
\]
For any current iterate $b$, the second-order integral expansion gives
\[
L(\beta)\leq Q(\beta\mid b):=
L(b)+\nabla L(b)^T(\beta-b)
+\tfrac18(\beta-b)^TX^TX(\beta-b).
\]
The first-order condition for minimizing $Q$ is
\[
-X^T\{y-p(b)\}+\tfrac14X^TX(\beta-b)=0.
\]
Therefore,
\[
\boxed{\beta^{(k+1)}=\beta^{(k)}+
4(X^TX)^{-1}X^T(y-p^{(k)}).}
\]
The majorization property gives
\[
L(\beta^{(k+1)})\leq Q(\beta^{(k+1)}\mid\beta^{(k)})
\leq Q(\beta^{(k)}\mid\beta^{(k)})=L(\beta^{(k)}),
\]
so the log-likelihood is nondecreasing.
The two probit algorithms use the updates derived in Exercise 2.2.

\textbf{Step 4: Initialization and convergence.}
The two starts are
\[
\beta^{(0)}=0,\qquad
\beta^{(0)}=(X^TX)^{-1}X^Ty.
\]
The second is the least-squares fit to the binary response, used only as an
initial value. For each start and method, use either
\[
\|\beta^{(k)}-\beta^{(k-1)}\|_2\leq\epsilon
\]
or
\[
|K_n(\beta^{(k)})-K_n(\beta^{(k-1)})|\leq\epsilon,
\qquad K_n=\ell/n,
\]
with $\epsilon\in\{10^{-3},10^{-6}\}$. This produces
$4\times2\times2\times2=32$ runs. The first coefficient update counts as
iteration 1. The implementation uses full steps, without step-halving.

\textbf{Step 5: Computational results.}
The accompanying R program reports the stopping iteration, log-likelihood,
and all coefficient estimates for every run. The following tables are
included automatically after running the script in the document directory.
\IfFileExists{exercise_2_3_iterations.tex}{
\begin{table}[htbp]
\centering
\small
\caption{Exercise 2.3: convergence results.}
\begin{tabular}{llllrrr}
\toprule
Link & Method & Start & Rule & Tolerance & $k$ & $K_n$ \\
\midrule
logit & NR & zero & beta & 1e-03 &  5 & -0.51713100 \\
logit & NR & zero & beta & 1e-06 &  6 & -0.51713100 \\
logit & NR & zero & loglik & 1e-03 &  3 & -0.51713139 \\
logit & NR & zero & loglik & 1e-06 &  4 & -0.51713100 \\
logit & NR & LS & beta & 1e-03 &  5 & -0.51713100 \\
logit & NR & LS & beta & 1e-06 &  6 & -0.51713100 \\
logit & NR & LS & loglik & 1e-03 &  3 & -0.51713134 \\
logit & NR & LS & loglik & 1e-06 &  4 & -0.51713100 \\
logit & MM & zero & beta & 1e-03 & 15 & -0.51713102 \\
logit & MM & zero & beta & 1e-06 & 31 & -0.51713100 \\
logit & MM & zero & loglik & 1e-03 &  4 & -0.51747089 \\
logit & MM & zero & loglik & 1e-06 & 11 & -0.51713162 \\
logit & MM & LS & beta & 1e-03 & 15 & -0.51713102 \\
logit & MM & LS & beta & 1e-06 & 31 & -0.51713100 \\
logit & MM & LS & loglik & 1e-03 &  4 & -0.51746844 \\
logit & MM & LS & loglik & 1e-06 & 11 & -0.51713162 \\
probit & NR & zero & beta & 1e-03 &  5 & -0.51583557 \\
probit & NR & zero & beta & 1e-06 &  5 & -0.51583557 \\
probit & NR & zero & loglik & 1e-03 &  3 & -0.51583561 \\
probit & NR & zero & loglik & 1e-06 &  4 & -0.51583557 \\
probit & NR & LS & beta & 1e-03 &  5 & -0.51583557 \\
probit & NR & LS & beta & 1e-06 &  5 & -0.51583557 \\
probit & NR & LS & loglik & 1e-03 &  3 & -0.51583561 \\
probit & NR & LS & loglik & 1e-06 &  4 & -0.51583557 \\
probit & FS & zero & beta & 1e-03 &  5 & -0.51583557 \\
probit & FS & zero & beta & 1e-06 &  8 & -0.51583557 \\
probit & FS & zero & loglik & 1e-03 &  3 & -0.51583858 \\
probit & FS & zero & loglik & 1e-06 &  5 & -0.51583557 \\
probit & FS & LS & beta & 1e-03 &  5 & -0.51583557 \\
probit & FS & LS & beta & 1e-06 &  8 & -0.51583557 \\
probit & FS & LS & loglik & 1e-03 &  3 & -0.51583819 \\
probit & FS & LS & loglik & 1e-06 &  5 & -0.51583557 \\
\bottomrule
\end{tabular}
\end{table}

}{\par\emph{Iteration table pending execution of the accompanying R script.}\par}
\IfFileExists{exercise_2_3_coefficients.tex}{
\begin{landscape}
\scriptsize
\begin{longtable}{rlllllllllll}
\caption{Exercise 2.3: coefficient estimates. Run numbers follow the convergence table.}\\
\toprule
Run & Intercept & age & lwt & raceblack & raceother & smoke & ptdTRUE & ht & ui & ftv1 & ftv2+ \\
\midrule
\endfirsthead
\toprule
Run & Intercept & age & lwt & raceblack & raceother & smoke & ptdTRUE & ht & ui & ftv1 & ftv2+ \\
\midrule
\endhead
1 & 0.823019 & -0.037234 & -0.015653 & 1.192413 & 0.740685 & 0.755528 & 1.343763 & 1.913166 & 0.680195 & -0.436380 & 0.179009 \\
2 & 0.823019 & -0.037234 & -0.015653 & 1.192413 & 0.740685 & 0.755528 & 1.343763 & 1.913166 & 0.680195 & -0.436380 & 0.179009 \\
3 & 0.818550 & -0.037039 & -0.015621 & 1.189761 & 0.738288 & 0.753411 & 1.342133 & 1.910263 & 0.679602 & -0.435638 & 0.178018 \\
4 & 0.823013 & -0.037234 & -0.015653 & 1.192409 & 0.740681 & 0.755525 & 1.343761 & 1.913162 & 0.680195 & -0.436379 & 0.179007 \\
5 & 0.823019 & -0.037234 & -0.015653 & 1.192413 & 0.740685 & 0.755528 & 1.343763 & 1.913166 & 0.680195 & -0.436380 & 0.179009 \\
6 & 0.823019 & -0.037234 & -0.015653 & 1.192413 & 0.740685 & 0.755528 & 1.343763 & 1.913166 & 0.680195 & -0.436380 & 0.179009 \\
7 & 0.818896 & -0.037053 & -0.015624 & 1.189956 & 0.738454 & 0.753546 & 1.342280 & 1.910514 & 0.679653 & -0.435725 & 0.178083 \\
8 & 0.823014 & -0.037234 & -0.015653 & 1.192410 & 0.740682 & 0.755525 & 1.343761 & 1.913162 & 0.680195 & -0.436379 & 0.179007 \\
9 & 0.821999 & -0.037194 & -0.015646 & 1.191832 & 0.740188 & 0.755085 & 1.343414 & 1.912519 & 0.680050 & -0.436211 & 0.178788 \\
10 & 0.823018 & -0.037234 & -0.015653 & 1.192413 & 0.740684 & 0.755528 & 1.343763 & 1.913165 & 0.680195 & -0.436380 & 0.179008 \\
11 & 0.685019 & -0.032116 & -0.014594 & 1.114313 & 0.676060 & 0.697951 & 1.297345 & 1.823875 & 0.659727 & -0.413190 & 0.150250 \\
12 & 0.817093 & -0.036999 & -0.015610 & 1.189045 & 0.737809 & 0.752965 & 1.341738 & 1.909407 & 0.679352 & -0.435402 & 0.177731 \\
13 & 0.822003 & -0.037194 & -0.015646 & 1.191834 & 0.740189 & 0.755087 & 1.343415 & 1.912521 & 0.680051 & -0.436211 & 0.178789 \\
14 & 0.823018 & -0.037234 & -0.015653 & 1.192413 & 0.740684 & 0.755528 & 1.343763 & 1.913165 & 0.680195 & -0.436380 & 0.179008 \\
15 & 0.685583 & -0.032129 & -0.014599 & 1.114570 & 0.676270 & 0.698094 & 1.297581 & 1.824298 & 0.659803 & -0.413380 & 0.150315 \\
16 & 0.817112 & -0.037000 & -0.015610 & 1.189056 & 0.737818 & 0.752974 & 1.341745 & 1.909420 & 0.679355 & -0.435406 & 0.177735 \\
17 & 0.461595 & -0.023301 & -0.008925 & 0.704881 & 0.441686 & 0.464632 & 0.816932 & 1.132274 & 0.414497 & -0.279566 & 0.086411 \\
18 & 0.461595 & -0.023301 & -0.008925 & 0.704881 & 0.441686 & 0.464632 & 0.816932 & 1.132274 & 0.414497 & -0.279566 & 0.086411 \\
19 & 0.460815 & -0.023260 & -0.008921 & 0.704433 & 0.441247 & 0.464211 & 0.816653 & 1.131908 & 0.414412 & -0.279341 & 0.086254 \\
20 & 0.461595 & -0.023301 & -0.008925 & 0.704881 & 0.441686 & 0.464632 & 0.816932 & 1.132274 & 0.414497 & -0.279565 & 0.086411 \\
21 & 0.461595 & -0.023301 & -0.008925 & 0.704881 & 0.441686 & 0.464632 & 0.816932 & 1.132274 & 0.414497 & -0.279566 & 0.086411 \\
22 & 0.461595 & -0.023301 & -0.008925 & 0.704881 & 0.441686 & 0.464632 & 0.816932 & 1.132274 & 0.414497 & -0.279566 & 0.086411 \\
23 & 0.460794 & -0.023259 & -0.008921 & 0.704411 & 0.441229 & 0.464192 & 0.816651 & 1.131893 & 0.414405 & -0.279339 & 0.086249 \\
24 & 0.461595 & -0.023301 & -0.008925 & 0.704881 & 0.441686 & 0.464631 & 0.816932 & 1.132274 & 0.414497 & -0.279565 & 0.086411 \\
25 & 0.461569 & -0.023299 & -0.008925 & 0.704864 & 0.441669 & 0.464609 & 0.816917 & 1.132269 & 0.414490 & -0.279541 & 0.086415 \\
26 & 0.461595 & -0.023301 & -0.008925 & 0.704881 & 0.441686 & 0.464632 & 0.816932 & 1.132274 & 0.414497 & -0.279566 & 0.086411 \\
27 & 0.457051 & -0.022931 & -0.008923 & 0.701767 & 0.438372 & 0.460482 & 0.814461 & 1.131470 & 0.413343 & -0.275885 & 0.086905 \\
28 & 0.461569 & -0.023299 & -0.008925 & 0.704864 & 0.441669 & 0.464609 & 0.816917 & 1.132269 & 0.414490 & -0.279541 & 0.086415 \\
29 & 0.461571 & -0.023299 & -0.008925 & 0.704865 & 0.441670 & 0.464610 & 0.816918 & 1.132270 & 0.414491 & -0.279542 & 0.086415 \\
30 & 0.461595 & -0.023301 & -0.008925 & 0.704881 & 0.441686 & 0.464632 & 0.816932 & 1.132274 & 0.414497 & -0.279566 & 0.086411 \\
31 & 0.457303 & -0.022956 & -0.008923 & 0.701984 & 0.438611 & 0.460774 & 0.814629 & 1.131596 & 0.413413 & -0.276080 & 0.086880 \\
32 & 0.461571 & -0.023299 & -0.008925 & 0.704865 & 0.441670 & 0.464610 & 0.816918 & 1.132270 & 0.414491 & -0.279542 & 0.086415 \\
\bottomrule
\end{longtable}
\end{landscape}

}{\par\emph{The 32-run coefficient table is generated by the same script.}\par}

\textbf{Step 6: Interpretation.}
The logistic algorithms optimize the same likelihood, as do the two probit
algorithms, so sufficiently accurate fits should agree within each link.
MM uses a fixed upper bound on curvature and is generally slower near the
optimum than Newton's method. Newton and scoring differ for probit because
their curvature matrices differ. A tighter tolerance cannot give an earlier
first stopping iteration along the same deterministic sequence of iterates.
A least-squares start is not guaranteed to be better in every configuration.
Finally, near a nonsingular maximizer,
\[
\ell(\widehat\beta)-\ell(\beta)
\approx\tfrac12(\beta-\widehat\beta)^TJ(\widehat\beta)
(\beta-\widehat\beta),
\]
so small likelihood changes need not imply equally small coefficient changes.
These are theoretical expectations to compare with the generated output,
not claims about unexecuted iteration counts.
\end{solution}

\question Exercise 2.4. Reproduce the coefficient estimates, model-based
standard errors, null deviance, and residual deviance for both links.
\begin{solution}
\textbf{Step 1: Accurate coefficient fits.}
Use $\beta^{(0)}=0$ and stop when
\[
\|\beta^{(k)}-\beta^{(k-1)}\|_2\leq10^{-10}.
\]
Use Newton--Raphson for logit and Fisher scoring for probit.

\textbf{Step 2: Model-based standard errors.}
The model-based covariance estimate is
\[
\widehat V_{\mathrm{model}}=I(\widehat\beta)^{-1},\qquad
\widehat{\operatorname{SE}}(\widehat\beta_j)
=\sqrt{[\widehat V_{\mathrm{model}}]_{jj}}.
\]
For logit,
\[
I(\widehat\beta)=X^T\operatorname{diag}
\{\widehat p_i(1-\widehat p_i)\}X.
\]
For probit,
\[
I(\widehat\beta)=X^T\operatorname{diag}
\left\{\frac{\phi(\widehat\eta_i)^2}
{\widehat p_i(1-\widehat p_i)}\right\}X.
\]
Expected information is used to reproduce the model-based probit standard
errors from \texttt{glm}.

\textbf{Step 3: Deviances.}
For individual Bernoulli observations, the saturated log-likelihood is zero,
using $0\log0=0$. Thus
\[
D_{\mathrm{res}}=2\{\ell_{\mathrm{sat}}-\ell(\widehat\beta)\}
=-2\ell(\widehat\beta).
\]
The intercept-only MLE has constant probability $\widehat p_0=\overline y$,
regardless of link. Hence
\[
D_{\mathrm{null}}=-2\sum_i\{y_i\log\overline y
+(1-y_i)\log(1-\overline y)\}.
\]
The degrees of freedom are $n-1=188$ and $n-p=178$, respectively.

\textbf{Step 4: Reference values.}
The following are the values printed in the exercise, provided as benchmarks.
Small last-digit differences may occur with a different convergence rule.
\begin{center}
\small
\begin{tabular}{lrrrr}
\toprule
&\multicolumn{2}{c}{Logit}&\multicolumn{2}{c}{Probit}\\
Term&Estimate&SE&Estimate&SE\\
\midrule
Intercept&0.82302&1.24471&0.461585&0.722771\\
age&-0.03723&0.03870&-0.023300&0.022514\\
lwt&-0.01565&0.00708&-0.008925&0.004056\\
raceblack&1.19241&0.53597&0.704876&0.318383\\
raceother&0.74069&0.46174&0.441680&0.266840\\
smoke&0.75553&0.42502&0.464624&0.246865\\
ptdTRUE&1.34376&0.48062&0.816927&0.288368\\
ht&1.91317&0.72074&1.132272&0.426167\\
ui&0.68019&0.46434&0.414494&0.281152\\
ftv1&-0.43638&0.47939&-0.279556&0.279582\\
ftv2+&0.17901&0.45638&0.086412&0.269392\\
\bottomrule
\end{tabular}
\end{center}
\begin{center}
\begin{tabular}{lrr}
\toprule
Quantity&Logit&Probit\\
\midrule
Null deviance&234.67&234.67\\
Residual deviance&195.48&194.99\\
Null degrees of freedom&188&188\\
Residual degrees of freedom&178&178\\
\bottomrule
\end{tabular}
\end{center}
The custom numerical estimates and standard errors are included with the
robust-SE tables in Exercise 2.5 after running the R code.
\end{solution}

\question Exercise 2.5. Compute model-robust standard errors for the
logistic and probit coefficient estimates.
\begin{solution}
\textbf{Step 1: Individual scores and asymptotic covariance.}
Let $s_i(\beta)=\partial\ell_i(\beta)/\partial\beta=x_iu_i(\beta)$.
The likelihood estimator solves $\sum_i s_i(\widehat\beta)=0$.
Under possible misspecification, define $\beta^\dagger$ by
$E\{s_i(\beta^\dagger)\}=0$. With suitable regularity conditions,
\[
\sqrt n(\widehat\beta-\beta^\dagger)
\overset{d}{\longrightarrow}N(0,A^{-1}BA^{-T}),
\]
where
\[
A=E\left\{-\frac{\partial s_i(\beta^\dagger)}{\partial\beta^T}\right\},
\qquad B=E\{s_i(\beta^\dagger)s_i(\beta^\dagger)^T\}.
\]
This follows by Taylor expanding the estimating equation around
$\beta^\dagger$, applying a law of large numbers to the score derivative,
and a central limit theorem to $n^{-1/2}\sum_i s_i(\beta^\dagger)$.

\textbf{Step 2: Empirical sandwich.}
Use the sum-based matrices
\[
J=-\sum_i\frac{\partial s_i(\widehat\beta)}{\partial\beta^T},
\qquad M=\sum_i s_i(\widehat\beta)s_i(\widehat\beta)^T.
\]
The HC0 covariance estimate is
\[
\boxed{\widehat V_{\mathrm{HC0}}=J^{-1}MJ^{-T},\qquad
\widehat{\operatorname{SE}}_{\mathrm{HC0}}(\widehat\beta_j)
=\sqrt{[\widehat V_{\mathrm{HC0}}]_{jj}}.}
\]
There is no additional factor of $n$ because $J$ and $M$ are sums.
Equivalently, with $\widehat A=J/n$ and $\widehat B=M/n$,
$\widehat V=n^{-1}\widehat A^{-1}\widehat B\widehat A^{-T}$.

\textbf{Step 3: Logistic case.}
Since $s_i=x_i(y_i-p_i)$,
\[
J_L=X^T\operatorname{diag}\{\widehat p_i(1-\widehat p_i)\}X,
\qquad
M_L=X^T\operatorname{diag}\{(y_i-\widehat p_i)^2\}X.
\]
Therefore,
\[
\boxed{\widehat V_L=J_L^{-1}M_LJ_L^{-1}.}
\]

\textbf{Step 4: Probit case.}
Here
\[
\widehat u_i=\frac{\phi(\widehat\eta_i)(y_i-\widehat p_i)}
{\widehat p_i(1-\widehat p_i)},\qquad
M_P=X^T\operatorname{diag}(\widehat u_i^2)X.
\]
Using the observed weights from Exercise 2.2,
\[
J_P=X^T\operatorname{diag}(\widehat w_i^{\mathrm{obs}})X,
\qquad
\boxed{\widehat V_P=J_P^{-1}M_PJ_P^{-1}.}
\]
For the model-robust probit sandwich, the bread is the empirical observed
score derivative, not the model-expected information. Those population
matrices need not agree under misspecification.

\textbf{Step 5: Implementation and output.}
The R function \texttt{inference} constructs a matrix whose $i$th row is
$s_i(\widehat\beta)^T$, obtains $M$ by its cross-product, and combines it
with the observed information. It reports both HC0 and the optional correction
\[
\widehat V_{\mathrm{HC1}}=\frac{n}{n-p}\widehat V_{\mathrm{HC0}}.
\]
HC0 is the direct sandwich answer; HC1 is labeled separately.
\IfFileExists{inference_logit.tex}{\input{inference_logit.tex}}{
\par\emph{Logit numerical inference table pending execution of the R script.}\par}
\IfFileExists{inference_probit.tex}{\input{inference_probit.tex}}{
\par\emph{Probit numerical inference table pending execution of the R script.}\par}

Robust standard errors leave the coefficient estimates unchanged and need
not exceed model-based standard errors. Under misspecification, they support
inference for the pseudo-true coefficient, not an automatically causal
parameter. The calculation assumes independent observations; clustering
would require an appropriately clustered meat matrix.
\end{solution}

\end{questions}
\end{document}
